\documentclass[10pt]{article}
\usepackage[margin=2.2cm,landscape]{geometry}
\input{glyphtounicode}\pdfgentounicode=1
\usepackage[T1]{fontenc}\usepackage{lmodern}
\usepackage{amsmath,amssymb,booktabs,tabularx,array,graphicx}
\renewcommand{\thetable}{S\arabic{table}}
\newcommand{\nVerified}{40}
\title{Supplementary material\\[4pt]\large Procedure Selection versus Workflow Structuring in Local RAG for Institutional SOP Question Answering: An Ablation Study}
\date{}
\begin{document}
\maketitle
\begin{table}[!htbp]
\centering
\caption{Complete evaluation results across all 10 systems on \nVerified{} verified benchmark questions, reporting mean [95\% bootstrap confidence intervals].}
\label{tab:results-full}
\scriptsize
\setlength{\tabcolsep}{2pt}
\resizebox{\textwidth}{!}{%
\begin{tabular}{lccccccc}
\toprule
System & KFR $\uparrow$ & Coverage $\uparrow$ & TIR $\downarrow$ & AMR $\downarrow$ & Refusal acc. $\uparrow$ & False ref. $\downarrow$ & Prompt tok. \\
\midrule
Naive RAG & 0.712 [0.571, 0.840] & 0.727 [0.455, 1.000] & 0.000 [0.000, 0.000] & 0.000 [0.000, 0.000] & 0.600 [0.200, 1.000] & 0.000 [0.000, 0.000] & 916 \\
GraphRAG (entity) & 0.698 [0.560, 0.829] & 0.687 [0.415, 0.920] & 0.083 [0.000, 0.250] & 0.000 [0.000, 0.000] & 0.600 [0.200, 1.000] & 0.000 [0.000, 0.000] & 949 \\
\textbf{Workflow GraphRAG} & 0.955 [0.917, 0.990] & 1.000 [1.000, 1.000] & 0.000 [0.000, 0.000] & 0.000 [0.000, 0.000] & 1.000 [1.000, 1.000] & 0.000 [0.000, 0.000] & 1525 \\
\textit{Selection-Matched Text RAG (ablation)} & 0.962 [0.914, 1.000] & 0.975 [0.926, 1.000] & 0.000 [0.000, 0.000] & 0.000 [0.000, 0.000] & 1.000 [1.000, 1.000] & 0.000 [0.000, 0.000] & 1479 \\
\midrule
Hybrid RAG & 0.683 [0.536, 0.819] & 0.712 [0.439, 0.970] & 0.000 [0.000, 0.000] & 0.000 [0.000, 0.000] & 0.600 [0.200, 1.000] & 0.000 [0.000, 0.000] & 882 \\
Agentic RAG & 0.712 [0.571, 0.840] & 0.702 [0.430, 0.950] & 0.000 [0.000, 0.000] & 0.000 [0.000, 0.000] & 0.600 [0.200, 1.000] & 0.000 [0.000, 0.000] & 1815 \\
Corrective RAG & 0.569 [0.414, 0.719] & 0.636 [0.364, 0.909] & 0.000 [0.000, 0.000] & 0.000 [0.000, 0.000] & 1.000 [1.000, 1.000] & 0.371 [0.229, 0.543] & 725 \\
Multimodal-layout RAG & 0.702 [0.562, 0.831] & 0.702 [0.430, 0.950] & 0.000 [0.000, 0.000] & 0.000 [0.000, 0.000] & 0.600 [0.200, 1.000] & 0.000 [0.000, 0.000] & 994 \\
\midrule
LLM only & 0.113 [0.050, 0.183] & 0.105 [0.000, 0.209] & 0.333 [0.000, 0.667] & 0.444 [0.000, 1.000] & 0.000 [0.000, 0.000] & 0.000 [0.000, 0.000] & 107 \\
Full context & 0.962 [0.914, 1.000] & 0.944 [0.871, 1.000] & 0.027 [0.000, 0.082] & 0.000 [0.000, 0.000] & 1.000 [1.000, 1.000] & 0.000 [0.000, 0.000] & 2722 \\
\bottomrule
\end{tabular}
}
\end{table}

\begin{table}[!htbp]
\centering
\caption{The 17 hand-curated relational triples of the custom entity-graph GraphRAG baseline, as defined in \texttt{src/engine.py}. Entities and relations are kept in their original Indonesian form.}
\label{tab:triples}
\scriptsize
\setlength{\tabcolsep}{4pt}
\begin{tabular}{rlll}
\toprule
\# & Head entity & Tail entity & Relation \\
\midrule
1 & cuti akademik & izin cuti & prosedur pengajuan \\
2 & cuti akademik & aktif kembali & alur lanjutan \\
3 & cuti akademik & ketua program studi & rekomendasi \& persetujuan \\
4 & cuti akademik & dekan & penerbitan SK izin \\
5 & cuti akademik & spp/ukt & syarat bebas tunggakan \\
6 & legalisir & ijazah & dokumen objek \\
7 & legalisir & transkrip & dokumen objek \\
8 & legalisir & subbagian akademik & verifikasi \& validasi \\
9 & beasiswa & rekomendasi beasiswa & dokumen pengajuan \\
10 & beasiswa & wakil dekan i & pengesahan surat \\
11 & ukt & keterlambatan pembayaran & permohonan dispensasi \\
12 & ukt & bagian keuangan fsm & verifikasi pembayaran \\
13 & irs & pengisian irs & registrasi akademik per semester \\
14 & irs & dosen wali & bimbingan \& persetujuan online \\
15 & irs & siap undip & portal sistem informasi \\
16 & organisasi mahasiswa & proposal kegiatan & pengajuan persetujuan \\
17 & organisasi mahasiswa & wakil dekan i & persetujuan kegiatan \\
\bottomrule
\end{tabular}
\end{table}

\begin{table}[!htbp]
\centering
\caption{One benchmark question per category from \texttt{sop\_benchmark\_v1.json}, with English glosses.}
\label{tab:examples}
\scriptsize
\setlength{\tabcolsep}{3pt}
\begin{tabularx}{\textwidth}{l l >{\raggedright\arraybackslash}X}
\toprule
ID & Category & Question (Indonesian) / English gloss \\
\midrule
CUTI-P1 & Procedural & \textit{Apa saja tahapan mengajukan cuti akademik di FSM dari awal sampai surat izin keluar, dan siapa yang mengerjakan tiap tahap?} \newline What are the steps for applying for academic leave at FSM, from the start until the permit letter is issued, and who performs each step? \\
CUTI-A1 & Actor & \textit{Setelah form cuti saya ditandatangani Kaprodi, form itu diserahkan ke siapa dan diproses oleh bagian apa?} \newline After the Head of Study Programme signs my leave form, to whom is it submitted and which unit processes it? \\
CUTI-D1 & Documents & \textit{Dokumen apa saja yang harus dilampirkan saat mengajukan cuti akademik?} \newline Which documents must be attached when applying for academic leave? \\
CUTI-T1 & Time/location & \textit{Berapa lama maksimal proses izin cuti akademik, dan di mana saya mengambil Surat Izin Dekan?} \newline What is the maximum processing time for an academic leave permit, and where do I collect the Dean's permit letter? \\
X-1 & Cross-SOP & \textit{Apa perbedaan dokumen yang dibawa untuk mengajukan cuti akademik dengan untuk aktif kembali setelah cuti?} \newline How do the documents required for academic leave differ from those required for reactivation after leave? \\
NA-1 & Unanswerable & \textit{Berapa biaya legalisir ijazah per lembar di FSM?} \newline How much does diploma legalisation cost per sheet at FSM? \\
\bottomrule
\end{tabularx}
\end{table}
\end{document}
